\documentclass[11pt,a4paper]{article}
\usepackage[margin=25mm]{geometry}
\usepackage{amsmath,amssymb,amsthm,booktabs,hyperref,array}
\setlength{\emergencystretch}{2em}
\newtheorem{proposition}{Proposition}
\title{Wilson trialities: current corrected source contract}
\author{Human Compression project: L023 complete corrected contract}\date{}
\begin{document}\maketitle
The fixed original is Dimitri Leemans and Klara Stokes,
\emph{Incidence geometries with trialities coming from maps with Wilson
trialities}, arXiv:2208.08215v2, 9 September 2022, thirteen physical pages.
This is a complete corrected principal submission for distinct assembly
review, not a claim that all literal original assertions are true. All five
complete written components are acquired, not treated as free citations:
\begin{center}
\begin{tabular}{l>{\raggedright\arraybackslash}p{9cm}}
4 pages&\emph{Cubic field automorphisms and Wilson trialities}: root's
exact native edition preserved in the distinct cubic-generation review.\\
3 pages&\emph{Exhausting the cubic matrix equations}: root's frozen
polynomial proof and complete required finite table program.\\
1 page&\emph{A boundary case for the rank-four map construction}:
root's digonal counterexample and its distinct review.\\
5 pages&\emph{Flag completion and Wilson symmetry in the four-type
construction}: new uniform proofs and exact small required certificate,
with root's distinct bounded review.\\
3 pages&\emph{Completing the odd projective-linear triality geometries}:
the uniform all-flags proof, including its precise signed-vector scope.\\
\end{tabular}
\end{center}
The preliminary status sentences in these immutable native documents
are historical. Current acceptance, authorship and limits are supplied
by the exact review bindings, not by silent edits to those documents.

\section*{Principal results and limits}
Source Theorem 3's odd-characteristic existence is proved by the complete
subfield-counting argument: every odd prime power $q$ admits a generating
marked triple for $\operatorname{PSL}_2(q^3)$ with Wilson class III.
The even-characteristic part is explicitly retained from Jones--Poulton
[source reference 18, Theorem 2.3], as in the original; we do not claim a
new proof of that classification/construction. The usual marked-map
correspondence is retained. It is not an assertion that every abstract
marked quotient is a faithful vertex--edge--face incidence map.

The polynomial component exhausts source equations (1)--(3), proves the
primitive-eigenvalue version of Lemma 5 and the characteristic-dependent
periods, and computes every one of the twelve rows of Table 1. All nine
odd rows agree. The three even rows are corrected to $(24,24)$,
$(216,144)$ and $(960,1920)$ for $q=4,8,16$ respectively. The additional
solutions have $bc=0$; this identifies the exact discrepancy, not the
source authors' unpublished computational procedure. A primitive root
in equation (4) is necessary for the stated large-order proof. The
source's experimental remark concerning other orders is not an exhaustive
new theorem and is not promoted to one here.
If ``root of unity'' in the fourth equation is read without primitivity,
there is an immediate literal counterexample to Lemma 5: in characteristic
two take $R_0=R_2=I$ and $\xi=1$. All four displayed equations hold,
but the projective group generated with $R_1=J$ has order two. This
does not dispute the primitive-root interpretation used for the corrected
finite columns.

Construction 6 always defines its literal four-type coset pregeometry.
The unconditional geometry claim, Lemma 7, is false even for a faithful
Class III map: the complete product example has group order 16128 and a
pairwise incident face--vertex--Petrie triple with no completing edge.
The corrected local theorem proves two free $G$-orbits of chambers
when the three nonedge parabolics are proper; an odd dihedron refutes
unrestricted Lemma 8. Digonal and dihedral examples refute unrestricted
Remark 9 and the Class I, II and IV correlation-quotient assertions.

The Class III symmetry assertion has a uniform faithful-map replacement.
Local incidence valencies intrinsically fix the edge type, and restriction
to the faithful original truncation excludes every type transposition.
Thus the exact type image is $C_3$, even without assuming the marked
group finite. The conclusion is about a pregeometry. For it to be an
incidence geometry, the separate all-flags completion condition must
hold. The explicit Class III counterexample proves that this qualification
cannot be removed.

\paragraph{The principal infinite-family conclusion.}
The fifth component supplies the missing existential conclusion by a
uniform written argument. For every odd prime power $q$, the accepted
marking in $G=\operatorname{PSL}_2(q^3)$ has a three-dimensional adjoint
model. A signed-vector stabilizer and norm argument bounds every
face--Petrie residue by three vertices; three forced vertices and exactly
two edges attain this bound. Thus every flag completes. Individual
isotropic vectors are allowed, and injectivity is asserted only on a
fixed face-star, not on the entire vertex set. The full automorphism
group is $G$, the full correlation group is $G\rtimes C_3$, and there
are two chamber orbits. Increasing $q$ gives nonisomorphic geometries.
This proves the advertised infinite existence while leaving the false
unrestricted construction explicitly refuted.

The source's final broad question asks for infinitely many
\emph{flag-transitive} geometries with trialities and no dualities.
It is answered by the separately accepted later Leemans--Stokes--Tranchida
result (L015): for every odd $n>1$ divisible by three, $\operatorname{Sz}(2^n)$
is the full group of such a thin, residually connected rank-three
geometry, with correlation group $G\rtimes C_3$. The accepted generation
repair is part of that attributed later result. This is not a claim of
rank-four flag transitivity, not a premise of our odd-PSL completion,
and not new question credit. Its exact theorem and acceptance are bound
separately in the historical disposition record.

\section*{Borrowed structure and one printed parity correction}
The entire subgroup classification in source Theorem 1 is borrowed
Dickson--Vanden Cruyce material, not newly proved by this assembly.
The actual structural inputs used by the new proofs are the subgroup
types, canonical subfield normalizers, Klein normalizers and involution
centralizers, the standard projective semilinear automorphism group,
and finite-field matrix Hilbert 90. Their complete external proof and
software costs are unknown. The cubic proof supplies the new uniform
bound on obstructing subfields rather than relying on the source's
claim that all proper stable generated groups lie in the cubic subfield.

Source Lemma 2 remains valid in its surrounding odd-$q$ context.
The cube of an odd residue is the same residue modulo eight. Thus the
retained Klein-normalizer classification gives the same order (12 or 24)
in $\operatorname{PSL}_2(q)$ and $\operatorname{PSL}_2(q^3)$.
The former normalizer is contained in the latter; their equal orders
prove equality. We do not extend this assertion to even characteristic.

For precision, one literal parameter in Theorem 1(6), physical page 6,
is reversed. Let $p$ be odd, $k\mid r$. An element of
$\mathbb F_{p^k}^{*}$ is a square in $\mathbb F_{p^r}^{*}$ according
to its square character raised to
\[
       1+p^k+\cdots+p^{r-k}\equiv r/k\pmod2.
\]
Consequently
\[
 \left|\mathbb F_{p^k}^{*}\cap(\mathbb F_{p^r}^{*})^2\right|
 =\begin{cases}p^k-1,&r/k\text{ even},\\
                (p^k-1)/2,&r/k\text{ odd}.
   \end{cases}
\]
These are the possible Borel multipliers in the ambient projective
special linear group. The source assigns the denominator 2 to the
even extension and 1 to the odd extension, in the opposite order.
For example its $q=3$, $s=k=r=1$ entry permits $E_3:C_2$, which
would be a subgroup of order six in $\operatorname{PSL}_2(3)=A_4$.
There is no such subgroup: an index-two subgroup would yield a
nontrivial map $A_4\to C_2$, whereas $A_4$ is generated by its
elements of order three. This is an explicit ancillary correction,
not a change to any accepted component's hypotheses.

\section*{Complete acquisition and attribution}
The source existence result and matrix equations are credited to
Leemans and Stokes, the even family to Jones and Poulton, and the retained
classification and map correspondence to the named external inputs.
Root authored the four-page cubic, three-page polynomial and one-page
digonal components. The present author distinctly reviewed the cubic
and digonal proofs and authored the five-page symmetry/counterexample component and the new
three-page completion proof. Root supplied the symmetry component's separate distinct review and scrutinized the completion proof. The polynomial
component has its own distinct review. No author self-accepts a new proof.

The complete \texttt{polynomial\_table.g} is required for the twelve
finite rows; the complete \texttt{certificate.py} is required for the
displayed fixed $\operatorname{SL}_2(8)$ instance. Their recorded output
is regenerated evidence, not an additional mathematical database.
All operational drivers and independent optional diagnostics are
separately inventoried in full. No new Lean declaration, source-count
admission, new question credit or numerical factor is supplied. The source
mentions Magma but supplies no complete executable listing; its software
cost is unknown, not zero. Any later measurement must preserve the
user's symmetric treatment of executable verification.
\end{document}
